\section{Probability Distributions}


\subsection{}
\label{2.1}
Let $x$ be a variable such that
%
\begin{equation}
p(x | \mu) = \mu ^ x (1 - \mu) ^ {1 - x},
\end{equation}
%
where $x \in \{ 0, 1 \}$.
Then,
%
\begin{equation}
\sum_{x} p(x | \mu) = 1.
\end{equation}
%
By the definition,
%
\begin{equation}
\begin{aligned}
{\rm E} x &= \mu, \\
{\rm E} x ^ 2 &= \mu,
\end{aligned}
\end{equation}
%
Since
%
\begin{equation}
{\rm var} x = {\rm E} x ^ 2 - \left( {\rm E} x \right) ^ 2,
\end{equation}
%
we have
%
\begin{equation}
{\rm var} x = \mu (1 - \mu).
\end{equation}
%
By the definition,
%
\begin{equation}
{\rm H} (x) = - \sum_{x} p(x | \mu) \ln p(x | \mu).
\end{equation}
%
Therefore,
%
\begin{equation}
{\rm H} (x) = - \mu \ln \mu - (1 - \mu) \ln (1 - \mu).
\end{equation}
%


\subsection{}
\label{2.2}
Let $x$ be a variable such that
%
\begin{equation}
p(x | \mu) = \left( \frac{1 - \mu}{2} \right) ^ \frac{1 - x}{2} \left( \frac{1 + \mu}{2} \right) ^ \frac{1 + x}{2},
\end{equation}
%
where $x \in \{ -1, 1 \}$.
Then,
%
\begin{equation}
\sum_{x} p(x | \mu) = 1.
\end{equation}
%
By the definition,
%
\begin{equation}
\begin{aligned}
{\rm E} x &= \mu, \\
{\rm E} x ^ 2 &= 1,
\end{aligned}
\end{equation}
%
Since
%
\begin{equation}
{\rm var} x = {\rm E} x ^ 2 - \left( {\rm E} x \right) ^ 2,
\end{equation}
%
we have
%
\begin{equation}
{\rm var} x = 1 - \mu ^ 2.
\end{equation}
%
By the definition,
%
\begin{equation}
{\rm H} (x) = - \sum_{x} p(x | \mu) \ln p(x | \mu).
\end{equation}
%
Therefore,
%
\begin{equation}
{\rm H} (x) = - \frac{1 - \mu}{2} \ln \frac{1 - \mu}{2} - \frac{1 + \mu}{2} \ln \frac{1 + \mu}{2}.
\end{equation}
%


\subsection{}
\label{2.3}
By the definition,
%
\begin{equation}
\begin{aligned}
{N \choose m} &= \frac{N!}{m! (N - m)!}, \\
{N \choose m - 1} &= \frac{N!}{(m - 1)! (N - m + 1)!}
\end{aligned}
\end{equation}
%
Therefore,
%
\begin{equation}
{N \choose m} + {N \choose m - 1} = \frac{(N - m + 1) N! + m N!}{m! (N - m + 1)!}.
\end{equation}
%
By the definition, the right hand side can be written as
%
\begin{equation}
\frac{(N + 1)!}{m! (N + 1 - m)!} = {N + 1 \choose m}.
\end{equation}
%
Thus,
%
\begin{equation}
{N \choose m} + {N \choose m - 1} = {N + 1 \choose m}.
\end{equation}
%

Note that 
%
\begin{equation}
1 + x = \sum_{m = 0}^{1} {1 \choose m} x ^ m.
\end{equation}
%
Let us assume that 
%
\begin{equation}
(1 + x) ^ N = \sum_{m = 0}^{N} {N \choose m} x ^ m.
\end{equation}
%
Then,
%
\begin{equation}
(1 + x) ^ {N + 1} = \sum_{m = 0}^{N} {N \choose m} x ^ m + \sum_{m = 0}^{N} {N \choose m} x ^ {m + 1}.
\end{equation}
%
By the result above, the right hand side can be written as
%
\begin{equation}
\sum_{m = 0}^{N} {N \choose m} x ^ m + \sum_{m = 1}^{N + 1} {N \choose m - 1} x ^ m = 1 + x ^ {N + 1} + \sum_{m = 1}^{N} {N + 1 \choose m} x ^ m.
\end{equation}
%
Therefore,
%
\begin{equation}
(1 + x) ^ {N + 1} = \sum_{m = 0}^{N + 1} {N + 1 \choose m} x ^ m.
\end{equation}
%
Thus, the assumption is proved by induction on $N$.

Finally, let $m$ be a variable such that
%
\begin{equation}
p(m | \mu) = {N \choose m} \mu ^ m (1 - \mu) ^ {N - m}.
\end{equation}
%
Then
%
\begin{equation}
\sum_{m = 0}^{N} p(m | \mu) = \sum_{m = 0}^{N} {N \choose m} \mu ^ m (1 - \mu) ^ {N - m}.
\end{equation}
%
By the result above, the right hand side can be written as
%
\begin{equation}
(1 - \mu) ^ N \sum_{m = 0}^{N} {N \choose m} \left( \frac{\mu}{1 - \mu} \right) ^ m = (1 - \mu) ^ N \left( 1 + \frac{\mu}{1 - \mu} \right) ^ N.
\end{equation}
%
Therefore,
%
\begin{equation}
\sum_{m = 0}^{N} p(m | \mu) = 1.
\end{equation}
%


\subsection{}
\label{2.4}
Let $m$ be a variable such that
%
\begin{equation}
p(m | \mu) = {N \choose m} \mu ^ m (1 - \mu) ^ {N - m}.
\end{equation}
%
Then
%
\begin{equation}
{\rm E} m = \sum_{m = 0}^{N} m {N \choose m} \mu ^ m (1 - \mu) ^ {N - m}.
\end{equation}
%
Differentiating both sides of 
%
\begin{equation}
\sum_{m = 0}^{N} {N \choose m} \mu ^ m (1 - \mu) ^ {N - m}  = 1
\end{equation}
%
with respect to $\mu$ gives
%
\begin{equation}
\sum_{m = 0}^{N} m {N \choose m} \mu ^ {m - 1} (1 - \mu) ^ {N - m} - \sum_{m = 0}^{N} (N - m) {N \choose m} \mu ^ m (1 - \mu) ^ {N - m - 1} = 0.
\end{equation}
%
The first term of the left hand side can be written as $\frac{1}{\mu} {\rm E} m$.
Since
%
\begin{equation}
(N - m) {N \choose m} = N {N - 1 \choose m},
\end{equation}
%
the second term of the left hand side can be written as
%
\begin{equation}
- N \sum_{m = 0}^{N - 1} {N - 1 \choose m} \mu ^ m (1 - \mu) ^ {N - m - 1} = - N.
\end{equation}
%
Therefore,
%
\begin{equation}
{\rm E} m = N \mu. 
\end{equation}
%

Differentiating both sides of 
%
\begin{equation}
\sum_{m = 0}^{N} {N \choose m} \mu ^ m (1 - \mu) ^ {N - m}  = 1
\end{equation}
%
twice with respect to $\mu$ gives
%
\begin{equation}
\begin{aligned}
&\sum_{m = 0}^{N} m (m - 1) {N \choose m} \mu ^ {m - 2} (1 - \mu) ^ {N - m} \\
&- 2 \sum_{m = 0}^{N} m (N - m) {N \choose m} \mu ^ {m - 1} (1 - \mu) ^ {N - m - 1} \\
&+ \sum_{m = 0}^{N} (N - m) (N - m - 1) {N \choose m} \mu ^ m (1 - \mu) ^ {N - m - 2} = 0.
\end{aligned}
\end{equation}
%
The first term of the left hand side can be written as $\frac{1}{\mu ^ 2} {\rm E} m (m - 1)$.
Since
%
\begin{equation}
\begin{aligned}
m (N - m) {N \choose m} &= N (N - 1) {N - 2 \choose m - 1}, \\
(N - m) (N - m - 1) {N \choose m} &= N (N - 1) {N - 2 \choose m},
\end{aligned}
\end{equation}
%
the second and third term of the left hand side can be written as
%
\begin{equation}
\begin{aligned}
- 2 N (N - 1) \sum_{m = 1}^{N - 1} {N - 2 \choose m - 1} \mu ^ {m - 1} (1 - \mu) ^ {N - m - 1} &= - 2 N (N - 1), \\
N (N - 1) \sum_{m = 0}^{N} {N - 2 \choose m} \mu ^ m (1 - \mu) ^ {N - m - 2} &= N (N - 1).
\end{aligned}
\end{equation}
%
Therefore,
%
\begin{equation}
{\rm E} m (m - 1) = N (N - 1) \mu ^ 2.
\end{equation}
%
Thus, since 
%
\begin{equation}
{\rm var} m = {\rm E} m (m - 1) + {\rm E} m - ( {\rm E} m ) ^ 2,
\end{equation}
%
we have
%
\begin{equation}
{\rm var} m = N \mu ( 1 - \mu).
\end{equation}
%


\subsection{}
\label{2.5}
By the definition,
%
\begin{equation}
\Gamma (a) \Gamma (b) = \int_{0}^{\infty} x ^ {a - 1} \exp (- x) dx \int_{0}^{\infty} y ^ {b - 1} \exp (- y) dy.
\end{equation}
%
By the transformation $t = x + y$, the right hand side can be written as
%
\begin{equation}
\begin{aligned}
&\int_{0}^{\infty} x ^ {a - 1} \left( \int_{x}^{\infty} (t - x) ^ {b - 1} \exp (- t) dt \right) dx \\
&=  \int_{0}^{\infty} \left( \int_{0}^{t} x ^ {a - 1} (t - x) ^ {b - 1} dx \right) \exp (- t) dt.
\end{aligned}
\end{equation}
%
By the transformation $x = t \mu$, the right hand side can be written as
%
\begin{equation}
\begin{aligned}
&\int_{0}^{\infty} \left( \int_{0}^{1} (t \mu) ^ {a - 1} t ^ {b - 1} (1 - \mu) ^ {b - 1} t d\mu \right) \exp (- t) dt \\
&= \int_{0}^{1} \mu ^ {a - 1} (1 - \mu) ^ {b - 1} d\mu \int_{0}^{\infty} t ^ {a + b - 1} \exp ( - t ) dt.
\end{aligned}
\end{equation}
%
By the definition, the second integral of the right hand side  can be written as $\Gamma(a + b)$.
%
Therefore,
%
\begin{equation}
\int_{0}^{1} \mu ^ {a - 1} (1 - \mu) ^ {b - 1} d\mu = \frac{\Gamma (a) \Gamma (b)}{\Gamma (a + b)}.
\end{equation}
%


\subsection{}
\label{2.6}
Let $\mu$ be a variable such that
%
\begin{equation}
p(\mu | a, b) = \frac{\Gamma(a + b)}{\Gamma(a) \Gamma(b)} \mu ^ {a - 1} (1 - \mu) ^ {b - 1}.
\end{equation}
%
Then
%
\begin{equation}
\begin{aligned}
{\rm E} \mu &= \frac{\Gamma(a + b)}{\Gamma(a) \Gamma(b)} \int_{0}^{1} \mu ^ a (1 - \mu) ^ {b - 1} d\mu, \\
{\rm E} \mu ^ 2 &= \frac{\Gamma(a + b)}{\Gamma(a) \Gamma(b)} \int_{0}^{1} \mu ^ {a + 1} (1 - \mu) ^ {b - 1} d\mu.
\end{aligned}
\end{equation}
%
Since
%
\begin{equation}
\begin{aligned}
\int_{0}^{1} \mu ^ a (1 - \mu) ^ {b - 1} d\mu &= \frac{\Gamma (a + 1) \Gamma (b)}{\Gamma (a + b + 1)}, \\
\int_{0}^{1} \mu ^ {a + 1} (1 - \mu) ^ {b - 1} d\mu &= \frac{\Gamma (a + 2) \Gamma (b)}{\Gamma (a + b + 2)},
\end{aligned}
\end{equation}
%
we have
%
\begin{equation}
\begin{aligned}
{\rm E} \mu &= \frac{a}{a + b}, \\
{\rm E} \mu ^ 2 &= \frac{a (a + 1)}{(a + b) (a + b + 1)}.
\end{aligned}
\end{equation}
%
Since
%
\begin{equation}
{\rm var} \mu = {\rm E} \mu ^ 2 - \left( {\rm E \mu} \right) ^ 2,
\end{equation}
%
we have
%
\begin{equation}
{\rm var} \mu = \frac{a b}{(a + b) ^ 2 (a + b + 1)}.
\end{equation}
%
Since
%
\begin{equation}
\frac{\partial}{\partial \mu} p( \mu | a, b) = \frac{\Gamma(a + b)}{\Gamma(a) \Gamma(b)} \mu ^ {a - 1} (1 - \mu) ^ {b - 1} \left( \frac{a - 1}{\mu} - \frac{b - 1}{1 - \mu} \right),
\end{equation}
%
we have
%
\begin{equation}
{\rm mode} \mu = \frac{a - 1}{a + b - 2}.
\end{equation}
%


\subsection{}
Let $m$ and $l$ be a variable such that
%
\begin{equation}
p(m, l | \mu) = {m + l \choose m} \mu ^ m (1 - \mu) ^ l,
\end{equation}
%
where
%
\begin{equation}
p(\mu | a, b) = \frac{\Gamma(a + b)}{\Gamma(a) \Gamma(b)} \mu ^ {a - 1} (1 - \mu) ^ {b - 1}.
\end{equation}
%
By \ref{2.6}, 
%
\begin{equation}
{\rm E} (\mu | a, b) = \frac{a}{a + b}.
\end{equation}
%
Note that
%
\begin{equation}
\mu_{\rm ML} = \frac{m}{m + l}.
\end{equation}
%
Since
%
\begin{equation}
p(\mu | m, l, a,  b) \propto p(m, l | \mu) p(\mu | a, b),
\end{equation}
%
we have
%
\begin{equation}
p(\mu | m, l, a,  b) = \frac{\Gamma(m + l + a + b)}{\Gamma (m + a) \Gamma (l + b)} \mu ^ {m + a - 1} (1 - \mu) ^ {l + b - 1}.
\end{equation}
%
Therefore, by \ref{2.6}, 
%
\begin{equation}
{\rm E} (\mu | m, l, a, b) = \frac{m + a}{m + l + a + b}.
\end{equation}
%
Thus,
%
\begin{equation}
{\rm E} (\mu | m, l, a, b) = \lambda \mu_{\rm ML} + (1 - \lambda) {\rm E} (\mu | a, b),
\end{equation}
%
where
%
\begin{equation}
\lambda = \frac{m + l}{m + l + a + b}.
\end{equation}
%


\subsection{(Incomplete)}
\label{2.8}
Let $x$ and $y$ be variables.
Then, by the definition,
%
\begin{equation}
{\rm E} x = \int x p(x) dx.
\end{equation}
%
The right hand side can be written as
%
\begin{equation}
\int x \left( \int p(x, y) dy \right) dx = \int \left( \int x p(x | y) dx \right) p(y) dy.
\end{equation}
%
Therefore,
%
\begin{equation}
{\rm E} x = {\rm E}_y \left( {\rm E}_x (x | y) \right).
\end{equation}
%

By the definition,
%
\begin{equation}
{\rm var} x = {\rm E} \left( x - {\rm E} x \right) ^ 2.
\end{equation}
%
By the result above, the right hand side can be written as
%
\begin{equation}
\begin{aligned}
&{\rm E}_y \left( {\rm E}_x \left( \left( x - {\rm E}_x (x | y) + {\rm E}_x (x | y) - {\rm E} x \right) ^ 2 | y \right) \right) \\
&= {\rm E}_y \left( {\rm E}_x \left( \left( x - {\rm E}_x (x | y) \right) ^ 2 | y \right) \right) \\
&+ 2 {\rm E}_y \left( {\rm E}_x \left( \left( x - {\rm E}_x (x | y) \right) \left( {\rm E}_x (x | y) - {\rm E} x \right) | y \right) \right) \\
&+ {\rm E}_y \left( {\rm E}_x \left( \left( {\rm E}_x (x | y) - {\rm E} x \right) ^ 2 | y \right) \right)
\end{aligned}
\end{equation}
%
Let us look at each term of the right hand side.
By the definition, the first term can be written as ${\rm E}_y \left( {\rm var}_x (x | y) \right)$.
The second term can be written as
%
\begin{equation}
2 {\rm E}_y \left( \left( {\rm E}_x (x | y) - {\rm E} x \right) {\rm E}_x \left( \left( x - {\rm E}_x (x | y) \right) | y \right) \right)
\end{equation}
%
By the result above, the third term can be written as
%
\begin{equation}
{\rm E}_y \left( {\rm E}_x (x | y) - {\rm E}_y \left( {\rm E}_x (x | y) \right) \right) ^ 2 = {\rm var}_y \left( {\rm E}_x (x | y) \right).
\end{equation}
%
Therefore,
%
\begin{equation}
{\rm var} x = {\rm E}_y \left( {\rm var}_x (x | y) \right) + {\rm var}_y \left( {\rm E}_x (x | y) \right).
\end{equation}
%


\subsection{(Incomplete)}
\label{2.9}
For a vector $\bm{\mu}$ in 2 dimensions, \ref{2.5} gives
%
\[
\int_{
\begin{subarray}{l}
\mu_1 + \mu_2 = 1 \\\\
\mu_1 \geq 0, \mu_2 \geq 0
\end{subarray}
} 
\mu_1 ^ {\alpha_1 - 1} \mu_2 ^ {\alpha_2 - 1} d\bm{\mu}
= \frac{\Gamma(\alpha_1) \Gamma(\alpha_2)}{\Gamma(\alpha_1 + \alpha_2)}.
\]
%
For a vector $\bm{\mu}$ in $M$ dimensions, let us assume that
%
\[ 
\int_{
\begin{subarray}{l}
\sum_{k = 1}^{M} \mu_k = 1 \\\\
\mu_k \geq 0
\end{subarray}
} 
\prod_{k = 1}^{M} \mu_k ^ {\alpha_k - 1} d\bm{\mu} 
= \frac{\prod_{k = 1}^{M} \Gamma(\alpha_k)}{\Gamma \left( \sum_{k = 1}^{M} \alpha_k \right)}.
\]
%
Then, for a vector $\bm{\mu}$ in $M + 1$ dimensions, 
%
\[
\int_{
\begin{subarray}{l}
\sum_{k = 1}^{M + 1} \mu_k = 1 \\\\
\mu_k \geq 0
\end{subarray}
}
\prod_{k = 1}^{M + 1} \mu_k ^ {\alpha_k - 1} d\bm{\mu} 
= \int_{0}^{1} \mu_{M + 1} ^ {\alpha_{M + 1} - 1} 
\left( \int_{
\begin{subarray}{l}
\sum_{k = 1}^{M} \mu'_k = 1 - \mu_{M + 1} \\\\
\mu'_k \geq 0
\end{subarray}
}
\prod_{k = 1}^{M} {\mu'_k} ^ {\alpha_k - 1} d\bm{\mu}' 
\right) d\mu_{M + 1}.
\]
%
where $\bm{\mu}'$ is the vector of the first $M$ elements of $\bm{\mu}$.
By the transformation 
%
\begin{equation}
\bm{\mu}'' = \frac{1}{1 - \mu_{M + 1}} \bm{\mu}',
\end{equation}
%
the right hand side can be written as
%
\[
\int_{0}^{1} \mu_{M + 1} ^ {\alpha_{M + 1} - 1} 
\left( \int_{
\begin{subarray}{l}
\sum_{k = 1}^{M} \mu''_k = 1 \\\\
\mu''_k \geq 0
\end{subarray}
} \left( \prod_{k = 1}^{M} \left( (1 - \mu_{M + 1}){\mu''_k} \right) ^ {\alpha_k - 1} \right) (1 - \mu_{M + 1}) ^ M d\bm{\mu}'' \right) d\mu_{M + 1},
\]
%
so that
\[
\int_{0}^{1} \mu_{M + 1} ^ {\alpha_{M + 1} - 1} (1 - \mu_{M + 1}) ^ {\Sigma_{k = 1}^{M} \alpha_k} 
\left( \int_{
\begin{subarray}{l}
\sum_{k = 1}^{M} \mu''_k = 1 \\\\
\mu''_k \geq 0
\end{subarray}
} \prod_{k = 1}^{M} {\mu''_k} ^ {\alpha_k - 1} d\bm{\mu}'' \right) d\mu_{M + 1}.
\]
By the assumption, it can be written as
%
\begin{equation}
\frac{\prod_{k = 1}^{M} \Gamma(\alpha_k)}{\Gamma \left( \sum_{k = 1}^{M} \alpha_k \right)} \frac{\Gamma(\alpha_{M + 1}) \Gamma \left( \sum_{k = 1}^{M} \alpha_k + 1 \right)}{\Gamma \left( \sum_{k = 1}^{M + 1} \alpha_k + 1 \right)}
= \frac{\sum_{k = 1}^{M} \alpha_k}{\sum_{k = 1}^{M + 1} \alpha_k} \frac{\prod_{k = 1}^{M + 1} \Gamma(\alpha_k)}{\Gamma \left( \sum_{k = 1}^{M + 1} \alpha_k \right)}.
\end{equation}
%
Therefore,
%
\begin{equation}
\int \prod_{k = 1}^{M + 1} \mu_k ^ {\alpha_k - 1} d\bm{\mu} = \frac{\prod_{k = 1}^{M + 1} \Gamma(\alpha_k)}{\Gamma \left( \sum_{k = 1}^{M + 1} \alpha_k \right)} ?
\end{equation}
%
Thus, the assumption is proved by induction on $M$.


\subsection{}
\label{2.10}
Let $\bm{\mu}$ be a variable such that
%
\begin{equation}
p(\bm{\mu} | \bm{\alpha}) = \frac{\Gamma \left( \sum_{k = 1}^{K} \alpha_k \right)}{\prod_{k = 1}^{K} \Gamma(\alpha_k)} \prod_{k = 1}^{K} \mu_k ^ {\alpha_k - 1}.
\end{equation}
%
Then
%
\begin{equation}
\begin{aligned}
{\rm E} \mu_j &= \int \mu_j p(\bm{\mu} | \bm{\alpha}) d\bm{\mu}, \\
{\rm E} \mu_j ^ 2 &= \int \mu_j ^ 2 p(\bm{\mu} | \bm{\alpha}) d\bm{\mu}, \\
{\rm E} \mu_j \mu_l &= \int \mu_j \mu_l p(\bm{\mu} | \bm{\alpha}) d\bm{\mu}.
\end{aligned}
\end{equation}
%
If $j \neq l$, then the right hand sides can be written as
%
\begin{equation}
\begin{aligned}
\frac{\Gamma \left( \sum_{k = 1}^{K} \alpha_k \right)}{\prod_{k = 1}^{K} \Gamma(\alpha_k)} \frac{\frac{\Gamma(\alpha_j + 1)}{\Gamma(\alpha_j)} \prod_{k = 1}^{K} \Gamma(\alpha_k)}{\Gamma \left( \sum_{k = 1}^{K} \alpha_k + 1 \right)} &= \frac{\alpha_j}{\sum_{k = 1}^{K} \alpha_k}, \\
\frac{\Gamma \left( \sum_{k = 1}^{K} \alpha_k \right)}{\prod_{k = 1}^{K} \Gamma(\alpha_k)} \frac{\frac{\Gamma(\alpha_j + 2)}{\Gamma(\alpha_j)} \prod_{k = 1}^{K} \Gamma(\alpha_k)}{\Gamma \left( \sum_{k = 1}^{K} \alpha_k + 2 \right)} &= \frac{\alpha_j (\alpha_j + 1)}{\sum_{k = 1}^{K} \alpha_k \left( \sum_{k = 1}^{K} \alpha_k + 1 \right)},  \\
\frac{\Gamma \left( \sum_{k = 1}^{K} \alpha_k \right)}{\prod_{k = 1}^{K} \Gamma(\alpha_k)} \frac{\frac{\Gamma(\alpha_j + 1) \Gamma(\alpha_l + 1)}{\Gamma(\alpha_j) \Gamma(\alpha_l)} \prod_{k = 1}^{K} \Gamma(\alpha_k)}{\Gamma \left( \sum_{k = 1}^{K} \alpha_k + 2 \right)} &= \frac{\alpha_j \alpha_l}{\sum_{k = 1}^{K} \alpha_k \left( \sum_{k = 1}^{K} \alpha_k + 1 \right)}.
\end{aligned}
\end{equation}
%
Therefore,
%
\begin{equation}
\begin{aligned}
{\rm E} \mu_j &= \frac{\alpha_j}{\sum_{k = 1}^{K} \alpha_k}. \\
{\rm E} \mu_j ^ 2 &= \frac{\alpha_j (\alpha_j + 1)}{\sum_{k = 1}^{K} \alpha_k \left( \sum_{k = 1}^{K} \alpha_k + 1 \right)}, \\
{\rm E} \mu_j \mu_l &= \frac{\alpha_j \alpha_l}{\sum_{k = 1}^{K} \alpha_k \left( \sum_{k = 1}^{K} \alpha_k + 1 \right)}.
\end{aligned}
\end{equation}
%
Since
%
\begin{equation}
\begin{aligned}
{\rm var} \mu_j &= {\rm E} \mu_j ^ 2 - \left( {\rm E} \mu_j \right) ^ 2, \\
{\rm cov} \left( \mu_j, \mu_l \right) &= {\rm E} \mu_j \mu_l - {\rm E} \mu_j {\rm E} \mu_l,
\end{aligned}
\end{equation}
%
we have
%
\begin{equation}
\begin{aligned}
{\rm var} \mu_j &= \frac{\alpha_j \left( \sum_{k = 1}^{K} \alpha_k - \alpha_j \right)}{\left( \sum_{k = 1}^{K} \alpha_k \right) ^ 2 \left( \sum_{k = 1}^{K} \alpha_k + 1 \right)}, \\
{\rm cov} \left( \mu_j, \mu_l \right) &= - \frac{\alpha_j \alpha_l}{\left( \sum_{k = 1}^{K} \alpha_k \right) ^ 2 \left( \sum_{k = 1}^{K} \alpha_k + 1 \right)}.
\end{aligned}
\end{equation}
%


\subsection{}
\label{2.11}
Let $\bm{\mu}$ be a variable such that
%
\begin{equation}
p(\bm{\mu} | \bm{\alpha}) = \frac{\Gamma \left( \sum_{k = 1}^{K} \alpha_k \right)}{\prod_{k = 1}^{K} \Gamma(\alpha_k)} \prod_{k = 1}^{K} \mu_k ^ {\alpha_k - 1}.
\end{equation}
%
Then
%
\begin{equation}
{\rm E} \ln \mu_j = \int \left( \ln \mu_j \right) p(\bm{\mu} | \bm{\alpha}) d\bm{\mu}.
\end{equation}
%
Since
%
\begin{equation}
\frac{\partial}{\partial \alpha_j} p(\bm{\mu} | \bm{\alpha}) = \left( \frac{\Gamma' \left( \sum_{k = 1}^{K} \alpha_k \right)}{\Gamma \left( \sum_{k = 1}^{K} \alpha_k \right)} - \frac{\Gamma'(\alpha_j)}{\Gamma(\alpha_j)} + \ln \mu_j \right) p(\bm{\mu} | \bm{\alpha}),
\end{equation}
%
we have
%
\begin{equation}
{\rm E} \ln \mu_j = \frac{\partial}{\partial \alpha_j} \int p(\bm{\mu} | \bm{\alpha}) d\bm{\mu} + \left( \psi (\alpha_j) - \psi \left( \sum_{k = 1}^{K} \alpha_k \right) \right) \int p(\bm{\mu} | \bm{\alpha}) d\bm{\mu},
\end{equation}
%
where
%
\begin{equation}
\psi(a) = \frac{d}{da} \ln \Gamma(a).
\end{equation}
%
Therefore,
%
\begin{equation}
{\rm E} \ln \mu_j = \psi(\alpha_j) - \psi \left( \sum_{k = 1}^{K} \alpha_k \right).
\end{equation}
%


\subsection{}
\label{2.12}
Let $x$ be a variable such that
%
\begin{equation}
p(x | a,  b) = \frac{1}{b - a},
\end{equation}
%
where $a < b$.
Then
%
\begin{equation}
\int_{a}^{b} p(x | a, b) dx = 1.
\end{equation}
%
Note that
%
\begin{equation}
\begin{aligned}
{\rm E} x &= \int_{a}^{b} x p(x | a, b) dx, \\
{\rm E} x ^ 2 &= \int_{a}^{b} x ^ 2 p(x | a, b) dx.
\end{aligned}
\end{equation}
%
The right hand sides can be written as
%
\begin{equation}
\begin{aligned}
\frac{1}{b - a} \int_{a}^{b} x dx &= \frac{1}{2} (a + b), \\
\frac{1}{b - a} \int_{a}^{b} x ^ 2 dx &= \frac{1}{3} \left( a ^ 2 + a b + b ^ 2 \right).
\end{aligned}
\end{equation}
%
Therefore,
%
\begin{equation}
\begin{aligned}
{\rm E} x &= \frac{1}{2} (a + b), \\
{\rm E} x ^ 2 &= \frac{1}{3} \left( a ^ 2 + a b + b ^ 2 \right).
\end{aligned}
\end{equation}
%
Since 
%
\begin{equation}
{\rm var} x = {\rm E} x ^ 2 - \left( {\rm E} x \right) ^ 2,
\end{equation}
%
we have
%
\begin{equation}
{\rm var} x = \frac{1}{12} (b - a) ^ 2.
\end{equation}
%


\subsection{}
\label{2.13}
Let $\mathbf{x}$ be a variable in $D$ dimensions and
%
\begin{equation}
\begin{aligned}
p(\mathbf{x}) &= \mathcal{N} (\mathbf{x} | \bm{\mu}, \bm{\Sigma}), \\
q(\mathbf{x}) &= \mathcal{N} (\mathbf{x} | \mathbf{m}, \mathbf{L}).
\end{aligned}
\end{equation}
%
Then, by the definition,
%
\begin{equation}
{\rm KL} (p || q) = - \int \mathcal{N} (\mathbf{x} | \bm{\mu}, \bm{\Sigma}) \ln \frac{\mathcal{N} (\mathbf{x} | \mathbf{m}, \mathbf{L})}{\mathcal{N} (\mathbf{x} | \bm{\mu}, \bm{\Sigma})} d\mathbf{x}.
\end{equation}
%
Since
%
\begin{equation}
\ln \frac{\mathcal{N} (\mathbf{x} | \mathbf{m}, \mathbf{L})}{\mathcal{N} (\mathbf{x} | \bm{\mu}, \bm{\Sigma})} = \ln \frac{(2 \pi) ^ {- \frac{D}{2}} \left( \abs{\det \mathbf{L}} \right) ^ {- \frac{1}{2}} \exp \left( - \frac{1}{2} (\mathbf{x} - \mathbf{m}) ^ \intercal \mathbf{L} ^ {- 1} (\mathbf{x} - \mathbf{m}) \right)}{(2 \pi) ^ {- \frac{D}{2}} \left( \abs{\det \bm{\Sigma}} \right) ^ {- \frac{1}{2}} \exp \left( - \frac{1}{2} (\mathbf{x} - \bm{\mu}) ^ \intercal \bm{\Sigma} ^ {- 1} (\mathbf{x} - \bm{\mu}) \right)},
\end{equation}
%
The right hand side can be written as
%
\begin{equation}
\begin{aligned}
&\frac{1}{2} \ln \abs{\frac{\det \mathbf{L}}{\det \bm{\Sigma}}} \int \mathcal{N} (\mathbf{x} | \bm{\mu}, \bm{\Sigma}) d\mathbf{x} \\
&+ \frac{1}{2} \int (\mathbf{x} - \mathbf{m}) ^ \intercal \mathbf{L} ^ {- 1} (\mathbf{x} - \mathbf{m}) \mathcal{N} (\mathbf{x} | \bm{\mu}, \bm{\Sigma}) d\mathbf{x} \\
&- \frac{1}{2} \int (\mathbf{x} - \bm{\mu}) ^ \intercal \bm{\Sigma} ^ {- 1} (\mathbf{x} - \bm{\mu}) \mathcal{N} (\mathbf{x} | \bm{\mu}, \bm{\Sigma}) d\mathbf{x}.
\end{aligned}
\end{equation}
%
Let us look at each term.
Since
%
\begin{equation}
\int \mathcal{N} (\mathbf{x} | \bm{\mu}, \bm{\Sigma}) d\mathbf{x} = 1,
\end{equation}
%
the first term can be written as $\frac{1}{2} \ln \abs{\frac{\det \mathbf{L}}{\det \bm{\Sigma}}}$.
Since
%
\begin{equation}
(\mathbf{x} - \mathbf{m}) ^ \intercal \mathbf{L} ^ {- 1} (\mathbf{x} - \mathbf{m}) = (\mathbf{x} - \bm{\mu} + \bm{\mu} - \mathbf{m}) ^ \intercal \mathbf{L} ^ {- 1} (\mathbf{x} - \bm{\mu} + \bm{\mu} - \mathbf{m}),
\end{equation}
%
the second term can be written as
%
\begin{equation}
\begin{aligned}
&\frac{1}{2} \int (\mathbf{x} - \bm{\mu}) ^ \intercal \mathbf{L} ^ {- 1} (\mathbf{x} - \bm{\mu}) \mathcal{N} (\mathbf{x} | \bm{\mu}, \bm{\Sigma}) d\mathbf{x} \\
&+ (\bm{\mu} - \mathbf{m}) ^ \intercal \mathbf{L} ^ {- 1} \int (\mathbf{x} - \bm{\mu}) \mathcal{N} (\mathbf{x} | \bm{\mu}, \bm{\Sigma}) d\mathbf{x} \\
&+ \frac{1}{2} (\bm{\mu} - \mathbf{m}) ^ \intercal \mathbf{L} ^ {- 1} (\bm{\mu} - \mathbf{m}) \int \mathcal{N} (\mathbf{x} | \bm{\mu}, \bm{\Sigma}) d\mathbf{x}.
\end{aligned}
\end{equation}
%
Since
%
\begin{equation}
\begin{aligned}
\int \mathcal{N} (\mathbf{x} | \bm{\mu}, \bm{\Sigma}) d\mathbf{x} &= 1, \\
\int \mathbf{x} \mathcal{N} (\mathbf{x} | \bm{\mu}, \bm{\Sigma}) d\mathbf{x} &= \bm{\mu}, \\
\int (\mathbf{x} - \bm{\mu}) (\mathbf{x} - \bm{\mu}) ^ \intercal \mathcal{N} (\mathbf{x} | \bm{\mu}, \bm{\Sigma}) d\mathbf{x} &= \bm{\Sigma},
\end{aligned}
\end{equation}
%
it can be written as
%
\begin{equation}
\frac{1}{2} \tr \left( \mathbf{L} ^ {- 1} \bm{\Sigma} \right)
+ \frac{1}{2} (\bm{\mu} - \mathbf{m}) ^ \intercal \mathbf{L} ^ {- 1} (\bm{\mu} -  \mathbf{m}).
\end{equation}
%
Since
%
\begin{equation}
\int (\mathbf{x} - \bm{\mu}) (\mathbf{x} - \bm{\mu}) ^ \intercal \mathcal{N} (\mathbf{x} | \bm{\mu}, \bm{\Sigma}) d\mathbf{x} = \bm{\Sigma},
\end{equation}
%
the third term can be written as
%
\begin{equation}
- \frac{1}{2} \tr \left( \bm{\Sigma} ^ {- 1} \bm{\Sigma} \right) = - \frac{D}{2}
\end{equation}
%
Therefore,
%
\begin{equation}
{\rm KL} (p || q) 
= \frac{1}{2} \left( \ln \abs{\frac{\det \mathbf{L}}{\det \bm{\Sigma}}} 
+ \tr \left( \mathbf{L} ^ {- 1} \bm{\Sigma} \right) 
+ (\bm{\mu} - \mathbf{m}) ^ \intercal \mathbf{L} ^ {- 1} (\bm{\mu} - \mathbf{m}) 
- D \right).
\end{equation}
%


\subsection{(Incomplete)}
\label{2.14}
Let
%
\begin{equation}
\begin{aligned}
L(p(\mathbf{x})) = &- \int p(\mathbf{x}) \ln p(\mathbf{x}) d\mathbf{x}
+ \lambda_1 \left( \int p(\mathbf{x}) d\mathbf{x} - 1 \right) \\
&+ \lambda_2 \left( \int \bm{\mu} ^ \intercal \mathbf{x} p(\mathbf{x}) d\mathbf{x} - \norm{\bm{\mu}} ^ 2 \right)
+ \lambda_3 \left( \int (\mathbf{x} - \bm{\mu}) ^ \intercal (\mathbf{x} - \bm{\mu}) p(\mathbf{x}) d\mathbf{x} - \tr \bm{\Sigma} \right).
\end{aligned}
\end{equation}
%
Then
%
\begin{equation}
\frac{\delta L(p(\mathbf{x}))}{\delta p(\mathbf{x})} 
= - \ln p(\mathbf{x}) - 1 + \lambda_1 + \lambda_2 \bm{\mu} ^ \intercal \mathbf{x} + \lambda_3 (\mathbf{x} - \bm{\mu}) ^ \intercal (\mathbf{x} - \bm{\mu}).
\end{equation}
%
Setting the left hand side to zero gives
%
\begin{equation}
p(\mathbf{x}) 
= \exp \left( - 1 + \lambda_1 + \lambda_2 \bm{\mu} ^ \intercal \mathbf{x} + \lambda_3 (\mathbf{x} - \bm{\mu}) ^ \intercal (\mathbf{x} - \bm{\mu}) \right),
\end{equation}
%
so that
%
\begin{equation}
p(\mathbf{x}) 
= \exp \left( 
- 1 
+ \lambda_1 
- \frac{\lambda_2 ^ 2}{4 \lambda_3 ^ 2} \norm{\bm{\mu}} ^ 2
+ \lambda_3 \left( \mathbf{x} - \left( 1 - \frac{\lambda_2}{2 \lambda_3} \right) \bm{\mu} \right) ^ \intercal \left( \mathbf{x} - \left( 1 - \frac{\lambda_2}{2 \lambda_3} \right) \bm{\mu} \right) 
\right).
\end{equation}
%
Substituting it to 
%
\begin{equation}
\begin{aligned}
\int p(\mathbf{x}) d\mathbf{x} &= 1, \\
\int \mathbf{x} p(\mathbf{x}) d\mathbf{x} &= \bm{\mu}, \\
\int (\mathbf{x} - \bm{\mu}) (\mathbf{x} - \bm{\mu}) ^ \intercal p(\mathbf{x}) d\mathbf{x} &= \bm{\Sigma},
\end{aligned}
\end{equation}
%
and the transformation 
%
\begin{equation}
\mathbf{y} = \mathbf{x} - \left( 1 - \frac{\lambda_2}{2 \lambda_3} \right) \bm{\mu}
\end{equation}
%
gives
%
\begin{equation}
\begin{aligned}
\exp \left( - 1 + \lambda_1 - \frac{\lambda_2 ^ 2}{4 \lambda_3 ^ 2} \norm{\bm{\mu}} ^ 2 \right) 
\int \exp \left( \lambda_3 \mathbf{y} ^ \intercal \mathbf{y} \right) d\mathbf{y} 
&= 1, \\
\exp \left( - 1 + \lambda_1 - \frac{\lambda_2 ^ 2}{4 \lambda_3 ^ 2} \norm{\bm{\mu}} ^ 2 \right) 
\int \left( \mathbf{y} + \left( 1 - \frac{\lambda_2}{2 \lambda_3} \right) \bm{\mu} \right) \exp \left( \lambda_3 \mathbf{y} ^ \intercal \mathbf{y} \right) d\mathbf{y} 
&= \bm{\mu}, \\
\exp \left( - 1 + \lambda_1 - \frac{\lambda_2 ^ 2}{4 \lambda_3 ^ 2} \norm{\bm{\mu}} ^ 2 \right) 
\int \left( \mathbf{y} - \frac{\lambda_2}{2 \lambda_3} \bm{\mu} \right) \left( \mathbf{y} - \frac{\lambda_2}{2 \lambda_3} \bm{\mu} \right) ^ \intercal \exp \left( \lambda_3 \mathbf{y} ^ \intercal \mathbf{y} \right) d\mathbf{y} 
&= \bm{\Sigma}.
\end{aligned}
\end{equation}
%
%By the transformation
%%
%\begin{equation}
%y = \sqrt{- \lambda_3} \left( x - \left( \mu - \frac{\lambda_2}{2 \lambda_3} \right) \right),
%\end{equation}
%%
%they can be written as
%%
%\begin{equation}
%\begin{aligned}
%\exp \left( - 1 + \lambda_1 - \frac{\lambda_2 ^ 2}{4 \lambda_3} \right) \int_{- \infty}^{\infty} \exp \left( - y ^ 2 \right) (- \lambda_3) ^ {- \frac{1}{2}} dy &= 1,\\
%\exp \left( - 1 + \lambda_1 - \frac{\lambda_2 ^ 2}{4 \lambda_3} \right) \int_{- \infty}^{\infty} \left( (- \lambda_3) ^ {- \frac{1}{2}} y + \mu - \frac{\lambda_2}{2 \lambda_3} \right) \exp \left( - y ^ 2 \right) (- \lambda_3) ^ {- \frac{1}{2}} dy &= \mu,\\
%\exp \left( - 1 + \lambda_1 - \frac{\lambda_2 ^ 2}{4 \lambda_3} \right) \int_{- \infty}^{\infty} \left( (- \lambda_3) ^ {- \frac{1}{2}} y - \frac{\lambda_2}{2 \lambda_3} \right) ^ 2 \exp \left( - y ^ 2 \right) (- \lambda_3) ^ {- \frac{1}{2}} dy &= \sigma ^ 2.\\
%\end{aligned}
%\end{equation}
%%
%Since
%%
%\begin{equation}
%\begin{aligned}
%\int_{- \infty}^{\infty} \exp \left( - y ^ 2 \right) dy &= \Gamma \left( \frac{1}{2} \right), \\
%\int_{- \infty}^{\infty} y \exp \left( - y ^ 2 \right) dy &= 0, \\
%\int_{- \infty}^{\infty} y ^ 2 \exp \left( - y ^ 2 \right) dy &= \Gamma \left( \frac{3}{2} \right),
%\end{aligned}
%\end{equation}
%%
%they can be written as
%%
%\begin{equation}
%\begin{aligned}
%\exp \left( - 1 + \lambda_1 - \frac{\lambda_2 ^ 2}{4 \lambda_3} \right) (- \lambda_3) ^ {- \frac{1}{2}} \Gamma \left( \frac{1}{2} \right) &= 1, \\
%\exp \left( - 1 + \lambda_1 - \frac{\lambda_2 ^ 2}{4 \lambda_3} \right) \left( \mu - \frac{\lambda_2}{2 \lambda_3} \right) (- \lambda_3) ^ {- \frac{1}{2}} \Gamma \left( \frac{1}{2} \right) &= \mu, \\
%\exp \left( - 1 + \lambda_1 - \frac{\lambda_2 ^ 2}{4 \lambda_3} \right) \left( (- \lambda_3) ^ {- \frac{3}{2}} \Gamma \left( \frac{3}{2} \right) + (- \lambda_3) ^ {- \frac{1}{2}} \frac{\lambda_2 ^ 2}{4 \lambda_3 ^ 2} \Gamma \left( \frac{1}{2} \right) \right) &= \sigma ^ 2.
%\end{aligned}
%\end{equation}
%%
%Thus,
%%
%\begin{equation}
%\begin{aligned}
%\lambda_1 &= 1 - \frac{1}{2} \ln \left( 2 \pi \sigma ^ 2 \right), \\
%\lambda_2 &= 0, \\
%\lambda_3 &= - \frac{1}{2 \sigma ^ 2},
%\end{aligned}
%\end{equation}
%%
%so that
%%
%\begin{equation}
%p(x) = \left( 2 \pi \sigma ^ 2 \right) ^ {- \frac{1}{2}} \exp \left( - \frac{1}{2 \sigma ^ 2} (x - \mu) ^ 2 \right).
%\end{equation}
%%


\subsection{}
\label{2.15}
Let $\mathbf{x}$ be a variable in $D$ dimensions such that
%
\begin{equation}
p(\mathbf{x}) = \mathcal{N} (\mathbf{x} | \bm{\mu}, \bm{\Sigma}).
\end{equation}
%
Then, by the definition,
%
\begin{equation}
{\rm H}(\mathbf{x}) = - \int \mathcal{N} \left( \mathbf{x} | \mathbf{\mu}, \mathbf{\Sigma} \right) \ln \mathcal{N} \left( \mathbf{x} | \mathbf{\mu}, \mathbf{\Sigma} \right) d\mathbf{x}.
\end{equation}
%
The right hand side can be written as
%
\begin{equation}
\begin{aligned}
&- \int \mathcal{N} \left( \mathbf{x} | \bm{\mu}, \bm{\Sigma} \right) \left( - \frac{D}{2} \ln (2 \pi) - \frac{1}{2} \ln \abs{\det \bm{\Sigma}} - \frac{1}{2} (\mathbf{x} - \bm{\mu}) ^ \intercal \bm{\Sigma} ^ {- 1} (\mathbf{x} - \bm{\mu}) \right) d\mathbf{x} \\
=& \left( \frac{D}{2} \ln (2 \pi) + \frac{1}{2} \ln \abs{\det \bm{\Sigma}} \right) \int \mathcal{N} \left( \mathbf{x} | \bm{\mu}, \bm{\Sigma} \right) d\mathbf{x} \\
&+ \frac{1}{2} \int (\mathbf{x} - \bm{\mu}) ^ \intercal \bm{\Sigma} ^ {- 1} (\mathbf{x} - \bm{\mu}) \mathcal{N} \left( \mathbf{x} | \bm{\mu}, \bm{\Sigma} \right) d\mathbf{x}.
\end{aligned}
\end{equation}
%
Since
%
\begin{equation}
\begin{aligned}
\int \mathcal{N} \left( \mathbf{x} | \bm{\mu}, \bm{\Sigma} \right) d\mathbf{x} &= 1, \\
\int (\mathbf{x} - \bm{\mu}) (\mathbf{x} - \bm{\mu}) ^ \intercal \mathcal{N} \left( \mathbf{x} | \bm{\mu}, \bm{\Sigma} \right) d\mathbf{x} &= \bm{\Sigma},
\end{aligned}
\end{equation}
%
the first and second term of the right hand side can be written as
%
\begin{equation}
\frac{D}{2} \ln (2 \pi) + \frac{1}{2} \ln \abs{\det \bm{\Sigma}}
\end{equation}
%
and
%
\begin{equation}
\frac{1}{2} \tr \left( \bm{\Sigma} ^ {- 1} \bm{\Sigma} \right) = \frac{D}{2}.
\end{equation}
%
Therefore,
%
\begin{equation}
{\rm H}(\mathbf{x}) = \frac{D}{2} \left( 1 + \ln (2 \pi) \right) + \frac{1}{2} \ln \abs{\det \bm{\Sigma}}.
\end{equation}
%




































